\chapter{Limites et allocation du capital}\label{ch:fm:limits-and-capital-allocation}

Une firme détient 213~millions de dollars de capital de risque face à six stratégies, l'expected shortfall de son P\&L annuel dans sa pire année sur cent.
Mesurées une à une, les stratégies exigeraient 394~millions de dollars, et trois d'entre elles rapporteraient moins que le coût du capital de la firme, fixé à 15~\%. Mesurées par
ce que chacune apporte aux pires années de la firme, deux le franchissent encore~; une troisième le dépasse presque trois fois~; et le
livre de suivi de tendance (trend following), qui exigerait à lui seul 69~millions de dollars, réduit le capital de la firme de 21~millions de dollars, car il gagne de l'argent les années où les autres en perdent. Qui supporte
la queue de distribution décide des stratégies que la firme conserve.

\section{À quoi sert un dispositif de limites}

Une limite de risque (livre~6, chapitre~29) plafonne une mesure d'une exposition. Une firme en possède des milliers~; ce qui les tient ensemble est un
dispositif.

\begin{definition}[Dispositif de limites]\label{def:fm:limits-and-capital-allocation:framework}
Le \emph{dispositif de limites}\index{dispositif de limites} (limit framework) d'une firme est l'ensemble de ses limites de risque et des règles qui les relient~: l'arbre des niveaux
(firme, branche, desk, stratégie, trader) auxquels les limites sont fixées, les mesures limitées à chaque niveau, la façon dont les limites d'un niveau
contraignent celles des niveaux inférieurs, qui peut fixer, modifier et relever temporairement chaque limite, et ce qui se passe lorsqu'une limite est dépassée.
\end{definition}

Un dispositif répartit l'\omterm{def:fm:the-head-of-desks-job:appetite}{appétit pour le risque} de la firme (chapitre~6) dans la hiérarchie de risque~; ses limites sont la \omterm{def:fm:the-head-of-desks-job:appetite}{déclaration d'appétit pour le risque} exprimée en chiffres. Il remplit
trois fonctions~: maintenir la firme dans son appétit, rendre visible le risque de chaque desk par rapport à un chiffre convenu, et provoquer un échange avant qu'une
limite ne soit franchie, et non après. Un dispositif qui ne fait qu'enregistrer les dépassements remplit la deuxième fonction seulement~; celui dont les limites sont relevées
dès qu'elles sont atteintes n'en remplit aucune (chapitre~12).

\begin{omfigure}
\begin{tikzpicture}[font=\footnotesize,>=stealth,box/.style={draw,rounded corners,minimum height=0.62cm,align=center,fill=omDef!10}]
\node[box,fill=omProp!12,minimum width=3.6cm] (f) at (0,3) {firme\\{\scriptsize VaR, stress, liquidité}};
\node[box,minimum width=3.0cm] (d1) at (-3.2,1.6) {desk\\{\scriptsize VaR, stress, perte}};
\node[box,minimum width=3.0cm] (d2) at (3.2,1.6) {desk\\{\scriptsize VaR, stress, perte}};
\node[box,fill=gray!10,minimum width=2.4cm] (s1) at (-4.6,0.2) {stratégie\\{\scriptsize VaR, perte, notionnel}};
\node[box,fill=gray!10,minimum width=2.4cm] (s2) at (-1.8,0.2) {stratégie\\{\scriptsize sensibilités}};
\node[box,fill=gray!10,minimum width=2.4cm] (s3) at (1.8,0.2) {stratégie\\{\scriptsize VaR, perte}};
\node[box,fill=gray!10,minimum width=2.4cm] (s4) at (4.6,0.2) {stratégie\\{\scriptsize concentration}};
\draw[->] (f) -- (d1); \draw[->] (f) -- (d2); \draw[->] (d1) -- (s1); \draw[->] (d1) -- (s2); \draw[->] (d2) -- (s3); \draw[->] (d2) -- (s4);
\node[draw,rounded corners,fill=omThm!10,align=left,font=\scriptsize,anchor=west] at (-6.3,-1.3)
  {chaque limite~: seuil indicatif (alerte) et seuil ferme (action)~; hausses temporaires avec échéance~;\\l'exposition d'un niveau inférieur compte aussi contre la limite de son parent};
\end{tikzpicture}

\omcaption{Un \omterm{def:fm:limits-and-capital-allocation:framework}{dispositif de limites} en arbre~: l'appétit de la firme est réparti vers les desks et les stratégies, chaque niveau limitant les mesures qui correspondent à
ses risques. Modèle~: \texttt{firm.limitalloc.Node}, \texttt{check}.}\label{fig:fm:limits-and-capital-allocation:tree}
\end{omfigure}

\section{Mesures et niveaux~: notionnel, sensibilités, VaR, stress, perte}

Aucune mesure isolée ne suffit, car chacune ne voit pas ce qu'une autre perçoit. Les limites de notionnel et de position sont simples et ne peuvent pas être contournées par un modèle, mais ne disent
rien du risque~; les limites de sensibilité (delta, vega, DV01, spread de crédit) sont précises pour les petits mouvements et aveugles aux grands~; la valeur en
risque et l'expected shortfall (livre~6, chapitre~21) agrègent les risques mais reposent sur un modèle et un historique~; les pertes de stress (chapitre~22 du même
livre) voient des scénarios que l'historique n'a pas connus~; les limites de perte (livre~11, chapitre~27) n'agissent qu'après coup. Un desk en porte quelques-unes de
chaque sorte, choisies selon ses risques.

\begin{definition}[Limite de concentration]\label{def:fm:limits-and-capital-allocation:conc}
Une \emph{limite de concentration}\index{limite de concentration} plafonne une exposition à un même nom, émetteur, secteur, contrepartie ou instrument, ou une position
exprimée en part du volume d'un marché ou de ses positions ouvertes, quelle que soit la diversification par ailleurs~: elle limite ce qu'on ne pourrait pas vendre ou couvrir rapidement.
\end{definition}

\begin{method}[Choisir les limites d'un desk]\label{met:fm:limits-and-capital-allocation:choose}
\begin{enumerate}
\item Lister les risques du desk et les mouvements qui lui feraient le plus de mal~; retenir une mesure qui voit chacun d'eux.
\item Fixer une limite de VaR ou d'ES pour l'agrégat, des limites de stress pour les scénarios que l'appétit de la firme désigne, des limites de sensibilité pour les grands
risques, une limite de perte, et des \omterm{def:fm:limits-and-capital-allocation:conc}{limites de concentration} pour ce qui ne peut pas être vendu en quelques jours (les cas du chapitre~28 le montrent).
\item Dimensionner chaque limite à partir de la part du desk dans l'appétit de la firme (chapitre~6), et vérifier que le desk peut mener son activité prévue avec 60--80~\%
de chacune.
\item Consigner qui peut modifier chaque limite et la façon dont les hausses temporaires expirent.
\end{enumerate}
\end{method}

\section{Utilisation, dépassements et hausses temporaires}

\begin{definition}[Limite ferme, limite indicative, utilisation de la limite, dépassement de limite]\label{def:fm:limits-and-capital-allocation:breach}
Une \emph{limite ferme}\index{limite ferme} (hard limit) est une limite qui ne doit pas être dépassée~: tout excès doit être résorbé immédiatement, en réduisant l'exposition, ou
soumis à une demande de hausse approuvée. Une \emph{limite indicative}\index{limite indicative} (soft limit) est un seuil inférieur dont le franchissement déclenche une alerte et un examen, et non une
action obligatoire. L'\emph{utilisation de la limite}\index{utilisation de la limite} est l'exposition exprimée en part de la limite. Un
\emph{dépassement de limite}\index{dépassement de limite} est une exposition supérieure à une limite ferme, consignée avec son ampleur, sa durée et sa résolution.
\end{definition}

La stratégie de vente de volatilité du chapitre a une limite de VaR quotidienne de 8, avec un seuil indicatif à 6. Après une hausse de la volatilité, sa VaR atteint 7.6~: un dépassement de la
\omterm{def:fm:limits-and-capital-allocation:breach}{limite indicative} et une utilisation de 95~\% de la \omterm{def:fm:limits-and-capital-allocation:breach}{limite ferme}. Le desk demande une hausse temporaire à 12 pendant seize jours~; l'utilisation tombe à 63~\% et le
position est maintenue (\cref{fig:fm:limits-and-capital-allocation:var}). Une hausse temporaire est une décision de prendre davantage de risque pour une durée déclarée~; celle qui
est renouvelée à chaque échéance devient une hausse permanente que personne n'a approuvée comme telle.

\begin{omfigure}
\begin{tikzpicture}
\begin{axis}[width=10.5cm,height=5.2cm,xlabel={\small jour},ylabel={\small VaR quotidienne (millions de dollars)},xmin=1,xmax=60,ymin=0,ymax=13,
  tick label style={font=\footnotesize},grid=major,grid style={gray!20},table/col sep=comma,
  legend cell align=left,legend style={font=\scriptsize,at={(0.5,-0.3)},anchor=north,/tikz/every even column/.append style={column sep=8pt}},legend columns=3]
\addplot[omDef,thick] table[x=day,y=var]{figdata/desk/07-limits-and-capital-allocation/varpath.csv};
\addplot[omProp,thick,dashed] table[x=day,y=soft]{figdata/desk/07-limits-and-capital-allocation/varpath.csv};
\addplot[omThm,thick,const plot] table[x=day,y=hard]{figdata/desk/07-limits-and-capital-allocation/varpath.csv};
\legend{VaR quotidienne,limite indicative,{limite ferme (hausse temporaire jours 25--40)}}
\end{axis}
\end{tikzpicture}

\omcaption{VaR quotidienne de la stratégie de vente de volatilité pendant une hausse de la volatilité, face à sa \omterm{def:fm:limits-and-capital-allocation:breach}{limite indicative} et à sa \omterm{def:fm:limits-and-capital-allocation:breach}{limite ferme}, relevée temporairement de 8
à 12 pour les jours 25 à 40. Trajectoire illustrative~; modèle \texttt{fm\_limits.var\_path}.}\label{fig:fm:limits-and-capital-allocation:var}
\end{omfigure}

La même vérification s'applique à chaque niveau. Dans la firme du chapitre, le premier jour du choc, lorsque la VaR de la vente de volatilité a doublé pour atteindre 6.4,
les VaR quotidiennes des six stratégies s'additionnent à 21.5~millions de dollars, face à une
limite de firme de 30~millions de dollars (utilisation de 72~\%), alors que la seule vente de volatilité est en dépassement de sa \omterm{def:fm:limits-and-capital-allocation:breach}{limite indicative}~: la firme reste largement dans son appétit et un desk ne l'est pas. Un dispositif qui ne contrôlerait que la firme ne verrait rien~; celui qui ne contrôlerait que les desks ne verrait jamais le jour où six desks à 70~\%
totalisent un dépassement au niveau de la firme.
totalisent un dépassement au niveau de la firme.

\section{Allocation du capital de risque}

\begin{definition}[Capital économique]\label{def:fm:limits-and-capital-allocation:ec}
Le \emph{capital économique}\index{capital économique} d'une firme est le capital qu'elle détient, selon sa propre mesure, face aux pertes inattendues~: typiquement
l'expected shortfall ou la valeur en risque de son P\&L sur un an à un niveau de confiance élevé, calculé sur son propre modèle de ses risques.
\end{definition}

La firme du chapitre (illustrative, en millions de dollars par an) gère six stratégies~: stat arb, suivi de tendance, credit carry, vente de volatilité, macro et market
making, avec un P\&L espéré de 11, 12, 8, 12, 6 et 10 et des volatilités de 25, 30, 20, 22, 20 et 8, corrélées à 0.2 sauf le suivi de tendance. Une année sur
vingt, un krach ajoute des pertes de 40, 60 et 110 au stat arb, au credit carry et à la vente de volatilité, et un gain de 30 au suivi de tendance. Son \omterm{def:fm:limits-and-capital-allocation:ec}{capital économique}, l'expected shortfall à 99~\%
du P\&L annuel sur 200\,000 années simulées, est de 212.8~millions de dollars pour un P\&L espéré de 59.1~millions de dollars.

Trois allocations répartissent ce capital (\cref{fig:fm:limits-and-capital-allocation:capital}). \emph{Autonome}~: l'expected shortfall propre à chaque stratégie~;
elles s'additionnent à 393.9~millions de dollars, bien plus que celui de la firme, parce que la firme est diversifiée. \emph{Euler} (livre~6, chapitre~20)~: chaque
stratégie, le P\&L espéré dans les années de queue de la firme, avec le signe changé~; ces montants s'additionnent exactement à celui de la firme. \emph{Incrémentale}~: l'expected
shortfall de la firme moins ce qu'il serait sans la stratégie.

\begin{proposition}[Contributions d'Euler à l'expected shortfall]\label{prop:fm:limits-and-capital-allocation:euler}
Soit $X=\sum_iX_i$ le P\&L de la firme et $\mathrm{ES}_\alpha(X)=-\E[X\mid X\le q_{1-\alpha}(X)]$. Pour une loi continue, la contribution d'Euler de la
stratégie $i$, $w_i\,\partial\mathrm{ES}_\alpha(\sum_jw_jX_j)/\partial w_i$ en $w=\mathbf 1$, vaut
$-\E[X_i\mid X\le q_{1-\alpha}(X)]$, et les contributions s'additionnent à $\mathrm{ES}_\alpha(X)$.
\end{proposition}

\begin{proof}[Preuve partielle]
L'expected shortfall est positivement homogène de degré un en les échelles $w$, donc le théorème d'Euler donne $\sum_iw_i\,\partial_i\mathrm{ES}=\mathrm{ES}$.
La dérivée de l'espérance conditionnelle à la queue par rapport à $w_i$ est l'espérance conditionnelle à la queue de $X_i$, car la dérivée du quantile
n'y contribue pas à la frontière (la densité de $X$ au quantile multipliée par une variation de longueur nulle de l'ensemble de queue)~; l'argument complet, avec les
conditions de régularité, est celui de Tasche (livre~6, chapitre~20). Sur des scénarios, la contribution est la moyenne de $-X_i$ sur les
$(1-\alpha)n$ pires scénarios de la firme, et la somme est exactement l'expected shortfall sur scénarios de la firme.
\end{proof}

\begin{omfigure}
\begin{tikzpicture}
\begin{axis}[width=11.5cm,height=5.4cm,ybar,bar width=7pt,ymin=-40,ymax=130,ylabel={\small capital (millions de \$)},
  xtick={0,1,2,3,4,5},xticklabels={{stat arb},suivi de tendance,{credit carry},{vente de volatilité},macro,{market making}},
  xticklabel style={font=\scriptsize,text width=1.5cm,align=center},tick label style={font=\footnotesize},table/col sep=comma,
  enlarge x limits=0.1,grid=major,grid style={gray!20},legend cell align=left,
  legend style={font=\scriptsize,at={(0.5,-0.3)},anchor=north,/tikz/every even column/.append style={column sep=8pt}},legend columns=3]
\addplot[fill=gray!40,draw=gray] table[x=k,y=standalone]{figdata/desk/07-limits-and-capital-allocation/capital.csv};
\addplot[fill=omDef!55,draw=omDef] table[x=k,y=euler]{figdata/desk/07-limits-and-capital-allocation/capital.csv};
\addplot[fill=omThm!50,draw=omThm] table[x=k,y=incremental]{figdata/desk/07-limits-and-capital-allocation/capital.csv};
\draw[black] (axis cs:-0.6,0) -- (axis cs:5.6,0);
\legend{expected shortfall autonome,contribution d'Euler,incrémentale}
\end{axis}
\end{tikzpicture}

\omcaption{\omterm{def:fm:limits-and-capital-allocation:ec}{Capital économique} (expected shortfall à 99~\% du P\&L annuel) alloué à six stratégies de trois façons. Les allocations autonomes s'additionnent à
394~millions de dollars, les contributions d'Euler à celui de la firme, soit 213~millions de dollars~; la contribution du suivi de tendance est négative parce qu'il gagne dans les
pires années de la firme. Firme illustrative, 200\,000 années simulées. Données~: \texttt{fm\_limits.allocation}.}\label{fig:fm:limits-and-capital-allocation:capital}
\end{omfigure}

\omcode{code/firm/limitalloc/firm_limitalloc.py}{83}{102}{Allocations autonomes, d'Euler et incrémentales de l'expected shortfall sur scénarios,
et la rentabilité du capital et la charge en capital.}\label{lst:fm:limits-and-capital-allocation:alloc}

\section{Facturer le capital}

\begin{definition}[Rentabilité du capital ajustée du risque, charge en capital]\label{def:fm:limits-and-capital-allocation:raroc}
La \emph{rentabilité du capital ajustée du risque}\index{rentabilité du capital ajustée du risque} (RAROC) d'une stratégie est son P\&L espéré divisé par le
\omterm{def:fm:limits-and-capital-allocation:ec}{capital économique} qui lui est alloué. Sa \emph{charge en capital}\index{charge en capital} est le taux de rendement minimal exigé (livre~6, chapitre~19) multiplié par ce capital,
déduite de son P\&L lorsque sa performance est mesurée et que son enveloppe de rémunération est dimensionnée.
\end{definition}

Une stratégie crée de la valeur lorsque son RAROC dépasse le taux exigé, c'est-à-dire lorsque son P\&L espéré dépasse sa \omterm{def:fm:limits-and-capital-allocation:raroc}{charge en capital}
(\cref{fig:fm:limits-and-capital-allocation:va}). Avec un taux exigé de 15~\%~:

\begin{center}
\small
\begin{tabular}{@{}lrrrrr@{}}
\toprule
& espéré & \multicolumn{2}{c}{RAROC} & \multicolumn{2}{c}{valeur ajoutée}\\
stratégie & P\&L & autonome & Euler & autonome & Euler\\
\midrule
stat arb & 10.9 & 16.7~\% & 22.8~\% & 1.12 & 3.74\\
suivi de tendance & 12.0 & 17.4~\% & (capital négatif) & 1.66 & 15.24\\
credit carry & 8.0 & 10.5~\% & 12.6~\% & $-3.50$ & $-1.55$\\
vente de volatilité & 12.1 & 9.8~\% & 11.1~\% & $-6.34$ & $-4.20$\\
macro & 6.0 & 12.6~\% & 43.3~\% & $-1.17$ & 3.93\\
market making & 10.0 & 84.9~\% & (capital quasi nul) & 8.24 & 10.02\\
\bottomrule
\end{tabular}
\end{center}

\begin{omfigure}
\begin{tikzpicture}
\begin{axis}[width=11.5cm,height=5.2cm,ybar,bar width=9pt,ymin=-8,ymax=17,ylabel={\small valeur ajoutée (millions de \$ par an)},
  xtick={0,1,2,3,4,5},xticklabels={{stat arb},suivi de tendance,{credit carry},{vente de volatilité},macro,{market making}},
  xticklabel style={font=\scriptsize,text width=1.5cm,align=center},tick label style={font=\footnotesize},table/col sep=comma,
  enlarge x limits=0.1,grid=major,grid style={gray!20},legend cell align=left,
  legend style={font=\scriptsize,at={(0.5,-0.3)},anchor=north,/tikz/every even column/.append style={column sep=8pt}},legend columns=2]
\addplot[fill=gray!40,draw=gray] table[x=k,y=va_standalone]{figdata/desk/07-limits-and-capital-allocation/capital.csv};
\addplot[fill=omDef!55,draw=omDef] table[x=k,y=va_euler]{figdata/desk/07-limits-and-capital-allocation/capital.csv};
\draw[black] (axis cs:-0.6,0) -- (axis cs:5.6,0);
\legend{facturé sur le capital autonome,facturé sur le capital d'Euler}
\end{axis}
\end{tikzpicture}

\omcaption{P\&L espéré diminué d'une charge de 15~\% sur le capital alloué. Sur le capital autonome, trois stratégies détruisent de la valeur~; sur le capital d'Euler, deux le font,
credit carry et vente de volatilité, les deux qui perdent le plus dans les pires années de la firme~; macro et suivi de tendance changent de signe ou croissent. Données~:
\texttt{fm\_limits.value\_added}.}\label{fig:fm:limits-and-capital-allocation:va}
\end{omfigure}

La facturation autonome pénalise les diversificateurs et avantage les stratégies dont les pertes surviennent ensemble~; la facturation d'Euler est cohérente avec le
capital effectif de la firme et récompense ce qui réduit la queue de la firme. Elle a son propre piège~: une stratégie dont la contribution est proche de zéro ou négative a un RAROC indéfini ou
négatif et une \omterm{def:fm:limits-and-capital-allocation:raroc}{charge en capital} qui est un crédit, et un desk rémunéré sur cette base en voudra davantage que ce que la queue de la firme peut absorber (la contribution
change lorsque le portefeuille change). Les firmes plafonnent donc les crédits, appliquent un plancher, ou utilisent l'allocation incrémentale pour les décisions d'entrée ou
de sortie d'une activité, et l'allocation d'Euler pour la charge courante.

\begin{method}[Fixer une charge en capital]\label{met:fm:limits-and-capital-allocation:charge}
\begin{enumerate}
\item Calculer le \omterm{def:fm:limits-and-capital-allocation:ec}{capital économique} de la firme sur des scénarios qui incluent ses années de krach, et non seulement son historique récent.
\item L'allouer par contributions d'Euler pour la charge mensuelle~; recalculer les contributions lorsque les positions changent de manière significative.
\item Utiliser le capital incrémental pour les décisions d'ajout ou de fermeture d'une stratégie.
\item Facturer le capital au taux exigé dans le P\&L de chaque desk avant que son enveloppe de rémunération ne soit dimensionnée (chapitre~10), avec un plancher pour les stratégies dont
la contribution est proche de zéro.
\end{enumerate}
\end{method}

Une firme qui redimensionne ses stratégies pour maximiser le P\&L espéré diminué d'une charge de 15~\% sur son expected shortfall total, chaque échelle étant bornée entre zéro et
deux fois la taille actuelle (\texttt{firm.limitalloc.best\_scales}), porte quatre d'entre elles au plafond, réduit credit carry à 1.57 fois et vente de volatilité à
0.57 fois, et fait passer la valeur ajoutée de la firme de 27.2~millions de dollars à 59.0~millions de dollars par an. Le plafond fait l'essentiel de ce travail~; le résultat utile est
la direction, les deux stratégies les plus exposées à la queue se réduisent pendant que les autres croissent.

\section{Tutoriel~: six stratégies et une queue de distribution}\label{tut:fm:limits-and-capital-allocation:tut}

\begin{tutorial}
\textbf{Objectif.} Vérifier un \omterm{def:fm:limits-and-capital-allocation:framework}{dispositif de limites} sur une journée, allouer le \omterm{def:fm:limits-and-capital-allocation:ec}{capital économique} de la firme selon trois méthodes et le facturer. \textbf{État final~:} le tableau
des RAROC et la \cref{fig:fm:limits-and-capital-allocation:va}.
\begin{tutsteps}
\item \textbf{Limites.} \texttt{fm\_limits.framework()} construit l'arbre~; \texttt{firm.limitalloc.check(framework(), daily\_var\_exposures(), 1)}
signale le \omterm{def:fm:limits-and-capital-allocation:breach}{dépassement de limite} indicative et le taux d'utilisation de 72~\% de la firme~; \texttt{framework(temp=\dots)} ajoute une hausse temporaire de 4 à la limite de VaR de la
vente de volatilité jusqu'au jour 10.
\item \textbf{Scénarios.} \texttt{fm\_limits.scenarios()} simule 200\,000 années des six stratégies, avec le krach.
\item \textbf{Allocation.} \texttt{allocation()} calcule les trois allocations et les RAROC (le \cref{lst:fm:limits-and-capital-allocation:alloc})~;
vérifier que les contributions d'Euler s'additionnent à l'ES de la firme (la \cref{prop:fm:limits-and-capital-allocation:euler}).
\item \textbf{Charge et redimensionnement.} \texttt{value\_added()} et \texttt{scaled()}.
\end{tutsteps}
\end{tutorial}

\textbf{Pour aller plus loin.} Supprimer le krach et observer l'allocation d'Euler du suivi de tendance devenir positive~; réduire de moitié la vente de volatilité et recalculer la
contribution de chaque stratégie.

\section{Construction~: limites et capital}\label{bld:fm:limits-and-capital-allocation:limitalloc}

\begin{build}
\bfield{Rôle} Le \omterm{def:fm:limits-and-capital-allocation:framework}{dispositif de limites} de la firme sous forme de données et son \omterm{def:fm:limits-and-capital-allocation:ec}{capital économique} traduit en prix~: quel desk respecte ses limites aujourd'hui, et ce que coûte chaque stratégie
en capital.
\bfield{Interface} \texttt{firm.limitalloc}~: \texttt{Limit(measure, soft, hard)}, \texttt{Node(name, limits, children, temp)}~;
\texttt{check(node, exposures, day)}~; \texttt{es(pnl, alpha)}, \texttt{allocate\_capital(scen, alpha)}~; \texttt{raroc}, \texttt{capital\_charge}~;
\texttt{best\_scales(scen, alpha, h\_k, lo, hi)}~; \texttt{to\_riskctl(node)}, la partie pré-trade, sous la forme de l'instantané
\texttt{firm.riskctl} du livre~11.
\bfield{Règles} Une hausse temporaire ne s'applique que jusqu'à son dernier jour~; l'exposition d'un parent est la somme de celles de ses enfants, sauf si elle est mesurée directement~;
les contributions d'Euler s'additionnent exactement à l'ES de la firme sur les mêmes scénarios.
\bfield{Tests d'acceptation} \texttt{code/firm/limitalloc/tests/}~: statuts de \omterm{def:fm:limits-and-capital-allocation:breach}{limite indicative} et ferme, dépassement d'un parent, hausse arrivant à échéance~; contributions d'Euler
s'additionnent, et les contributions autonomes dépassent l'ES de la firme~; les allocations incrémentales ne dépassent pas les allocations autonomes~; l'optimiseur préfère la meilleure
stratégie~; la forme de l'export.
\bfield{Pour aller plus loin} Mesures non additives au niveau parent (la VaR d'un parent calculée à partir du P\&L conjoint)~; un optimiseur avec les limites de la firme comme contraintes~;
un plancher sur la charge.
\end{build}

\begin{omsources}
\item E.~Zaik, J.~Walter, G.~Kelling et C.~James, «~RAROC at Bank of America: from theory to practice~», \emph{Journal of Applied Corporate
Finance} 9(2), 1996.
\item Financial Stability Board, \emph{Principles for an Effective Risk Appetite Framework}, 2013.
\item Le livre~6, chapitres 19--21 et 29 (taux de rendement minimal exigé, allocation d'Euler, VaR et ES, limites du moteur de risque).
\end{omsources}

\section{Exercices}

\begin{exercise}[$\star$]\label{exo:fm:limits-and-capital-allocation:1}
La VaR quotidienne d'une stratégie vaut 7.6 face à une \omterm{def:fm:limits-and-capital-allocation:breach}{limite indicative} de 6 et une \omterm{def:fm:limits-and-capital-allocation:breach}{limite ferme} de 8. Quels sont son taux d'utilisation et son statut~? Et si la
\omterm{def:fm:limits-and-capital-allocation:breach}{limite ferme} est relevée temporairement à 12~?
\end{exercise}

\begin{exercise}[$\star$]\label{exo:fm:limits-and-capital-allocation:2}
Calculer le RAROC de la stat arb sur capital autonome et sur capital d'Euler, ainsi que sa valeur ajoutée à un taux de rendement minimal exigé de 15~\% pour chacun.
\end{exercise}

\begin{exercise}[$\star$]\label{exo:fm:limits-and-capital-allocation:3}
Les ES autonomes des six stratégies s'additionnent à 393.9~millions de dollars, et celui de la firme vaut 212.8~millions de dollars. Quel est le bénéfice de diversification, en proportion
de la somme des ES autonomes~?
\end{exercise}

\begin{exercise}[$\star\star$]\label{exo:fm:limits-and-capital-allocation:4}
À quel taux de rendement minimal exigé la vente de volatilité atteindrait-elle son point d'équilibre sur son capital d'Euler~? Et le credit carry~?
\end{exercise}

\begin{exercise}[$\star\star$]\label{exo:fm:limits-and-capital-allocation:5}
Expliquer pourquoi la contribution d'Euler du suivi de tendance est négative, et pourquoi un desk facturé sur cette base pourrait prendre plus de risque que la firme ne le souhaite.
\end{exercise}

\begin{exercise}[$\star\star$]\label{exo:fm:limits-and-capital-allocation:6}
Six desks utilisent chacun 70~\% d'une limite de VaR de 8, et la limite de la firme est de 30. La firme dépasse-t-elle sa limite si les VaR des desks s'additionnent~? Et si elles sont corrélées à
0.2 et si chaque VaR est un multiple de sa volatilité~?
\end{exercise}

\begin{exercise}[$\star\star\star$]\label{exo:fm:limits-and-capital-allocation:7}
\emph{Programmation.} Relancer \texttt{fm\_limits.allocation} sans krach (\texttt{P\_CRASH = 0}). Quelles sont désormais les
contributions d'Euler du suivi de tendance et de la vente de volatilité~?
\end{exercise}

\begin{exercise}[$\star\star\star$]\label{exo:fm:limits-and-capital-allocation:8}
\emph{Trouver la faille.} «~Chaque desk respecte ses limites, donc la firme respecte son \omterm{def:fm:the-head-of-desks-job:appetite}{appétit pour le risque}.~»
\end{exercise}

\section{Problème~: qui paie pour la queue}

\begin{problem}[{Problème du week-end --- qui paie pour la queue}]\label{pb:fm:limits-and-capital-allocation:1}
Le directeur des risques de la firme doit décider comment facturer six stratégies pour 213~millions de dollars de \omterm{def:fm:limits-and-capital-allocation:ec}{capital économique}, et le desk de vente de volatilité vient de demander
une hausse temporaire de limite.

\medskip\noindent\textbf{Partie I --- Le dispositif.}
\begin{enumerate}
\item Définir un \omterm{def:fm:limits-and-capital-allocation:framework}{dispositif de limites}, ainsi que les \omterm{def:fm:limits-and-capital-allocation:breach}{limites fermes} et les \omterm{def:fm:limits-and-capital-allocation:breach}{limites indicatives}.
\item Donner le taux d'utilisation et le statut de la stratégie de vente de volatilité au pic, avant et après une hausse temporaire à 12.
\item Donner le taux d'utilisation de la VaR de la firme le jour du pic.
\item Pourquoi un \omterm{def:fm:limits-and-capital-allocation:framework}{dispositif de limites} doit-il vérifier à la fois la firme et chaque desk~?
\item Qu'est-ce qui transforme une hausse temporaire en une hausse permanente non approuvée~?
\end{enumerate}

\medskip\noindent\textbf{Partie II --- Capital.}
\begin{enumerate}[resume]
\item Définir le \omterm{def:fm:limits-and-capital-allocation:ec}{capital économique} et donner celui de la firme, avec son P\&L espéré.
\item Donner le capital autonome, le capital d'Euler et le capital incrémental de chaque stratégie.
\item Énoncer et démontrer la \cref{prop:fm:limits-and-capital-allocation:euler}~; vérifier que les contributions s'additionnent.
\item Pourquoi la somme des capitaux autonomes dépasse-t-elle le capital de la firme~?
\item Quelle stratégie voit son capital changer de signe entre l'allocation autonome et l'allocation d'Euler, et pourquoi~?
\end{enumerate}

\medskip\noindent\textbf{Partie III --- La charge.}
\begin{enumerate}[resume]
\item Définir le RAROC et la \omterm{def:fm:limits-and-capital-allocation:raroc}{charge en capital}.
\item Quelles stratégies détruisent de la valeur avec un taux de rendement minimal exigé de 15~\%, selon chaque allocation~?
\item Donner la valeur ajoutée de chaque stratégie selon les deux allocations.
\item Qu'est-ce qui ne va pas avec une contribution d'Euler proche de zéro ou négative, et que font les firmes à ce sujet~?
\item Quelle allocation doit décider de la fermeture d'une stratégie~?
\end{enumerate}

\medskip\noindent\textbf{Partie IV --- La décision.}
\begin{enumerate}[resume]
\item Donner les meilleures échelles dans $[0,2]$ ainsi que la valeur ajoutée avant et après.
\item Quelle partie de ce résultat faut-il retenir, et laquelle faut-il écarter~?
\item Le desk de vente de volatilité doit-il obtenir sa hausse temporaire~?
\item Énoncer le \emph{résultat nommé}~: les taux de rendement minimaux auxquels le carry de crédit et la vente de volatilité atteignent le seuil de rentabilité, sur capital autonome et sur capital d'Euler, ainsi que
le capital du suivi de tendance (trend following) sous les deux allocations.
\item En deux phrases, rédiger la recommandation du responsable des risques.
\end{enumerate}
\end{problem}

\section{Questions d'entretien}

\begin{interviewq}[$\star$ \iqroles{risque}]\label{iq:fm:limits-and-capital-allocation:1}
Quelle est la différence entre une \omterm{def:fm:limits-and-capital-allocation:breach}{limite indicative} (soft limit) et une \omterm{def:fm:limits-and-capital-allocation:breach}{limite ferme} (hard limit)~? Donnez un exemple de chacune sur un desk de trading.
\end{interviewq}

\begin{interviewq}[$\star$ \iqroles{trader, risque}]\label{iq:fm:limits-and-capital-allocation:2}
Pourquoi un desk se dote-t-il de plusieurs types de limites plutôt que d'une seule limite de VaR~?
\end{interviewq}

\begin{interviewq}[$\star\star$ \iqroles{risque, chercheur}]\label{iq:fm:limits-and-capital-allocation:3}
Trois stratégies ont chacune un expected shortfall (ES) autonome de 50, et celui de la firme est de 90. Comment répartiriez-vous les 90, et pourquoi pas 50 chacune~?
\end{interviewq}

\begin{interviewq}[$\star\star$ \iqroles{banque, risque}]\label{iq:fm:limits-and-capital-allocation:4}
Une stratégie rapporte 12 par an sur un capital d'Euler de 108. Avec un taux de rendement minimal exigé de 15~\%, crée-t-elle de la valeur~? Qu'est-ce qui modifierait votre réponse~?
\end{interviewq}

\begin{interviewq}[$\star\star$ \iqroles{chercheur}]\label{iq:fm:limits-and-capital-allocation:5}
Pourquoi la contribution d'Euler d'une stratégie à l'expected shortfall peut-elle être négative~? Est-ce une raison pour la faire croître sans limite~?
\end{interviewq}

\begin{interviewq}[$\star\star\star$ \iqroles{risque}]\label{iq:fm:limits-and-capital-allocation:6}
Comment estimeriez-vous la contribution de chaque stratégie à la queue de distribution de la firme lorsque celle-ci est gouvernée par des krachs rares que les cinq dernières
années ne contiennent pas~?
\end{interviewq}
